%=============================================================================!
% 16.E  EXERCISES                                                             !
%=============================================================================!
\unnumbered{Exercises}
\label{sec:ob-exercises}

The exercises run from questions of understanding, through calculations,
to short derivations marked `show that'. Full solutions are in
\hyperref[sec:ob-solutions]{`Solutions to the exercises'} at the end of
the book, and Notebook~16 checks every numerical answer.

\begin{exercise}[the volume of a shell]
\label{ex:ob-shell}
For bins of \SI{0.05}{\angstrom}, by what fraction does $4\pi r^2\,\delta r$
fall short of the exact volume of the shell from $r$ to $r + \delta r$,
at $r = 1$ and at \SI{5}{\angstrom}? Show that the fraction is close to
$\delta r/r$ when $\delta r \ll r$.
\end{exercise}

\begin{exercise}[the far value of g(r)]
\label{ex:ob-far}
What value does $g(r)$ approach far from any atom, with ASE's
normalisation, for 256 and for 2048 atoms placed independently at random?
Which normalisation would make it approach 1 for them, and why is the
difference unimportant for the coordination number of the first shell?
\end{exercise}

\begin{exercise}[the shells of a face-centred cubic crystal]
\label{ex:ob-fcc}
In the fcc lattice of cube side $a$, the atoms sit at the corners of the
cube and at the centres of its faces. Show that an atom's nearest
neighbours are $a/\sqrt2$ away and number 12, and that the next shell is
at $a$ and holds 6.
\end{exercise}

\begin{exercise}[structure without structure]
\label{ex:ob-meanfield}
Suppose $g(r) = 0$ for $r < \sigma$ and $g(r) = 1$ beyond. Show that the
Lennard-Jones energy per atom is then $-\frac{16}{9}\pi\rho\sigma^3
\varepsilon$, and evaluate it for liquid argon, $\rho\sigma^3 = 0.8$ and
$\varepsilon = \SI{0.01034}{\electronvolt}$. Compare it with
\SI{-56.16}{\milli\electronvolt}, the energy per atom that the run's $g(r)$
gives with the full Lennard-Jones potential (the run itself switches the
potential off between $2\sigma$ and $2.5\sigma$), and say what the
difference measures.
\end{exercise}

\begin{exercise}[random atoms and a perfect crystal]
\label{ex:ob-sq-limits}
Show from \cref{eq:ob-sq} that $S(\vb{G})$ averages to 1 for atoms placed
independently at random in the cell, at any allowed $\vb{G} \neq
\vb{0}$, and is $N$ for a perfect crystal at a vector of its reciprocal
lattice.
\end{exercise}

\begin{exercise}[compressibility from scattering]
\label{ex:ob-s0}
A liquid of number density \SI{0.03}{\per\angstrom\cubed} at
\SI{300}{\kelvin} has $\kappa_T = \SI{0.5}{\per\giga\pascal}$. What is
$S(0)$? Would a liquid twice as compressible scatter long waves more
strongly or less?
\end{exercise}

\begin{exercise}[flicker]
\label{ex:ob-flicker}
A bonded pair vibrates with a period of \SI{60}{\femto\second}, and the
far end of each vibration carries it beyond a single cutoff. How many
events per picosecond does the cutoff record for this one pair? With
$r_{\mathrm{on}}$ at the old cutoff, how wide must the band of hysteresis
be to record none, and how wide to record none wherever the cutoff falls
within the vibration?
\end{exercise}

\begin{exercise}[bonds as keys]
\label{ex:ob-keys}
Six atoms have the bonds 0-1, 1-2, 3-4 in one frame and 0-1, 2-3, 3-4,
4-5 in the next. Write both sets as keys $iN + j$ with $N = 6$, find the
bonds formed and broken by the merge of sorted lists, and find the
molecules of each frame by the search of the toolbox in \cref{sec:ob-bonds}.
\end{exercise}

\begin{exercise}[two graphs that look alike]
\label{ex:ob-hash}
Six atoms of one element form a triangular prism (two triangles, each
atom of one joined to the atom above it in the other); six others form
a hexagon with each atom also joined to the opposite one. Show that every
atom in both has three neighbours, and explain why every round of the
hash of \cref{sec:ob-bonds} then gives every atom of both the same
number. Are the two molecules the same?
\end{exercise}

\begin{exercise}[how many events]
\label{ex:ob-poisson}
A rate is wanted to 2\%. How many events must be counted? A run of
\SI{1}{\nano\second} of 100 molecules shows 40 dissociations: give the
rate per molecule with its error.
\end{exercise}

\begin{exercise}[a random walk in one dimension]
\label{ex:ob-walk}
An atom on a line hops by $+a$ or $-a$ with equal probability, at the
rate $k_{\mathrm h}$. Show that its MSD after a time $t$ is $k_{\mathrm
h}a^2t$ and that $D = k_{\mathrm h}a^2/2$.
\end{exercise}

\begin{exercise}[hops that keep their direction]
\label{ex:ob-correlated}
On a line, each step of length $a$ repeats the direction of the one
before with probability $\frac12(1 + c)$, so that consecutive steps have
the correlation $c$, and steps $k$ apart $c^k$. Show that for many steps
the MSD is $na^2(1 + c)/(1 - c)$, and so that $f = (1 + c)/(1 - c)$, the
statistical inefficiency of \cref{sec:cv-correlation}. What $c$ gives the
memory of direction, $4D/k_{\mathrm h}\ens{\ell^2} = 3.46$, of the model
surface at \SI{1}{\per\pico\second}?
\end{exercise}

\begin{exercise}[the end of the ballistic part]
\label{ex:ob-water}
For Chapter~11's water, $D = \SI{5.24e-4}{\angstrom\squared\per\femto
\second}$ at \SI{304}{\kelvin}. Estimate when a molecule's motion stops
being ballistic, taking its mass as \SI{18.015}{\amu}.
\end{exercise}

\begin{exercise}[an exponential memory]
\label{ex:ob-exp}
Show from \cref{eq:ob-msd-vacf} that $C_{vv} = C_0\mathrm{e}^{-t/\tau_{\mathrm
c}}$ gives $\text{MSD} = 2C_0\tau_{\mathrm c}[t - \tau_{\mathrm c}(1 -
\mathrm{e}^{-t/\tau_{\mathrm c}})]$, that this is $C_0t^2$ for $t \ll
\tau_{\mathrm c}$ and $2C_0\tau_{\mathrm c}t$ for $t \gg \tau_{\mathrm c}$,
and so that $D = C_0\tau_{\mathrm c}/3$.
\end{exercise}

\begin{exercise}[the viscosity on the back of an envelope]
\label{ex:ob-eta}
\Cref{tab:cv-observables} gives the shear stress of liquid argon a
spread of \SI{0.0121}{\giga\pascal} and an integrated correlation time of
\SI{0.18}{\pico\second}. Show that the Green-Kubo integral is $(V/\kB T)
\sigma_{xy}^2\tau_{\mathrm{int}}$, with $\sigma_{xy}$ the spread, and
evaluate it
for 256 atoms at \SI{0.02035}{\per\angstrom\cubed} and \SI{135}{\kelvin}.
\end{exercise}

\begin{exercise}[Nernst-Einstein for the salt]
\label{ex:ob-ne}
From the mean diffusion coefficients of \cref{sec:ob-charge}, $5.19$ and
$\SI{5.03e-4}{\angstrom\squared\per\femto\second}$ for 32 ions of each
kind in a cell of side \SI{13.4}{\angstrom} at \SI{1487}{\kelvin},
compute $\sigma_{\mathrm{NE}}$ in \si{\siemens\per\metre}.
\end{exercise}

\begin{exercise}[one number for layers]
\label{ex:ob-trace}
What would a single diffusion coefficient, a sixth of the slope of the
full MSD, report for the guests of \cref{sec:ob-layers}? Why does it
describe neither of their motions?
\end{exercise}

\begin{exercise}[Euler's formula at work]
\label{ex:ob-euler}
Show that $\lvert\mathrm{e}^{\mathrm{i}\theta}\rvert = 1$, and recover
$\cos(A + B)$ and $\sin(A + B)$ from $\mathrm{e}^{\mathrm{i}A}
\mathrm{e}^{\mathrm{i}B} = \mathrm{e}^{\mathrm{i}(A + B)}$.
\end{exercise}

\begin{exercise}[the transform of one cosine]
\label{ex:ob-dft-cos}
Show that the discrete transform of $x_n = \cos(2\pi mn/N_{\mathrm s})$,
for a whole number $0 < m < N_{\mathrm s}/2$, is $N_{\mathrm s}/2$ at $k =
m$ and at $k = N_{\mathrm s} - m$ and zero elsewhere.
\end{exercise}

\begin{exercise}[how fast is fast]
\label{ex:ob-fft-cost}
Show that $W(N_{\mathrm s}) = 2W(N_{\mathrm s}/2) + N_{\mathrm s}$ with $W(1)
= 0$ gives $W(N_{\mathrm s}) = N_{\mathrm s}\log_2N_{\mathrm s}$ for
$N_{\mathrm s}$ a power of 2, and find the ratio of $N_{\mathrm s}^2$ to
$N_{\mathrm s}\log_2N_{\mathrm s}$ for $N_{\mathrm s} = 2^{20}$.
\end{exercise}

\begin{exercise}[a stride that hides a bond]
\label{ex:ob-fold}
The chain's stretch at \SI{18.25}{\tera\hertz} is seen in frames every
\SI{35}{\femto\second}. Where does it appear? What is the longest stride
at which it appears where it is?
\end{exercise}

\begin{exercise}[friction against resolution]
\label{ex:ob-width}
A spectrum is computed from a run under Langevin friction of
\SI{1}{\per\pico\second} with correlations taken to \SI{10}{\pico\second}.
Which widens its peaks more, the friction or the window?
\end{exercise}

\begin{exercise}[diffusion in a spectrum]
\label{ex:ob-vdos-zero}
What is $\mathcal{D}(0)$, in \si{\per\tera\hertz}, for the oxygen atoms of
Chapter~11's water, $m = \SI{15.999}{\amu}$, with the $D$ of
\cref{ex:ob-water}?
\end{exercise}

\begin{exercise}[where the free energy is lowest]
\label{ex:ob-fd}
Show that for $U \approx \frac12kr_{\mathrm h}^2$ the free energy along the distance
$r_{\mathrm h}$ from a hollow, $F(r_{\mathrm h}) = -\kB T\ln[r_{\mathrm h}\,\mathrm{e}^{-\beta kr_{\mathrm h}^2/2}] +
\text{constant}$, is least at $r_{\mathrm h} = \sqrt{\kB T/k}$, and evaluate it at
\SI{300}{\kelvin} for $k = \SI{0.9785}{\electronvolt\per\angstrom\squared}$.
\end{exercise}

\begin{exercise}[crossings]
\label{ex:ob-crossings}
A run crosses a barrier between two wells $n$ times. Argue that the free
energy difference between the wells is known to about $2\kB T/\sqrt n$,
and find how many crossings give it to $0.1\,\kB T$.
\end{exercise}
